\chapter{Ambivalence Principle: conjugate orientations \texorpdfstring{\(\pi/\varphi^2\)}{π/φ²} and \texorpdfstring{\(\varphi/\pi^2\)}{φ/π²}, sheet involution, and Euclidean realization of the equivalence principle.}
\label{chap:83-16}

\section{Compatible sections and measurement theory}
\label{p12-83-c16-s1:chap:medida}

The dimensional notion of measurement is formulated on inverse systems of
compatible states. Before assigning a number to a state, one specifies its refinement
history, the coupling with the apparatus, and the map that produces the
observable. This sequence makes it possible to separate the existence of a real value from
the physical operation by which that value is recorded.

\subsection{The number as a compatible section}

Let \((\mathcal S_n,r_{n+1,n})_{n\geq0}\) be an inverse system of
surviving states. A compatible section is a sequence
\(x=(x_n)_{n\geq0}\) such that
\[
 r_{n+1,n}(x_{n+1})=x_n
 \qquad(n\geq0).
\]
His space is the inverse limit.
\[
 \mathcal X_\infty=\varprojlim_n\mathcal S_n.
\]
When each \(x_n\) determines an interval \(I_n(x_n)\subset\mathbb R\), the
closures \(\overline{I_n(x_n)}\) are nonempty compact sets, are nested, and
their diameters tend to zero; the nested-interval theorem guarantees
that their intersection contains a unique point. The Archimedean evaluation is
then defined by
\[
 \mathcal P_{\rm arq}(x)
 =\text{the unique element of }
   \bigcap_{n\geq0}\overline{I_n(x_n)}.
\]
The HMT number is the compatible section
\(x\in\mathcal X_\infty\). The reader \(L\) is additional structure, and the
coupling \((x,L)\) constitutes its readout presentation or measurement. Two
sections may have the same real evaluation and differ in their quotient,
orientation, or boundary coordinates; the Archimedean projection does not identify
those genealogies by definition.

\subsection{Compatibility between ternary cylinders and thousandths}

A ternary word of length \(L\), whose integer value is \(N\), determines
the cylinder
\[
 I_{N,L}=\left[\frac{N}{3^L},\frac{N+1}{3^L}\right).
\]
A word of \(m\) blocks in base \(1000\), whose integer value is \(D\),
determines
\[
 J_{D,m}=\left[\frac{D}{1000^m},\frac{D+1}{1000^m}\right).
\]

\begin{lemma}[Exact compatibility of cylinders]
The cylinders \(I_{N,L}\) and \(J_{D,m}\) have a nonempty intersection if and
only if
\[
 N\,1000^m<(D+1)3^L,
 \qquad
 D\,3^L<(N+1)1000^m.
\]
\end{lemma}
\begin{proof}
Two half-open intervals \([a,b)\) and \([c,d)\) intersect exactly
when \(a<d\) and \(c<b\). Substituting the four endpoints and multiplying
by the positive integer \(3^L1000^m\) yields the inequalities. The
decision therefore reduces to integer arithmetic.
\end{proof}

The prolongation of a word refines its cylinder. A compatible chain whose closures are compact, nonempty, and nested, with diameter tending to zero, determines a unique real evaluation. The use of closures is essential: nested half-open intervals may have empty intersection when the limit coincides with an excluded endpoint. The complete state preserves the words, transformations, and carries that produce the evaluation.

\subsection{Reversible coupling between system and apparatus}

Let \(X\) be the state space of the system, \(M\) the state space of the
apparatus, and \(m_0\in M\) its initial preparation. A coupling is a
map
\[
 U:X\times M\longrightarrow X\times M.
\]
When \(U\) is bijective, the joint state preserves the information.
An evaluation map \(q_M:M\to Y\) may be noninjective and publish only an
observable \(y\in Y\).

\begin{definition}[Measurement]
An HMT measurement is the triple
\[
 \mathfrak M=(U,q_M,\mathcal C),
\]
where \(U\) couples system and apparatus, \(q_M\) evaluates the
apparatus state, and \(\mathcal C\subseteq X\times M\) is the compatibility condition. Its output
on \(x\) is
\[
 \operatorname{Med}_{\mathfrak M}(x)
 =q_M\!\left(\operatorname{pr}_M U(x,m_0)\right),
\]
provided that the coupled pair satisfies \(\mathcal C\).
\end{definition}

The definition establishes a precise composition of three operations:
coupling, restriction to the compatible domain, and evaluation of the observable.

\begin{theorem}[Localization of Information Loss]
Let
\[
 \Omega_M:=q_M\circ\operatorname{pr}_M:X\times M\longrightarrow Y.
\]
If \(U\) is bijective and the complete final state is preserved, the classes of
inputs indistinguishable under the published observable are exactly the
preimages under \(U\) of the fibres of \(\Omega_M\), unless it is further
composed with a subsequent elimination or reinitialization of the state.
\end{theorem}
\begin{proof}
The bijectivity of \(U\) permits \((x,m_0)\) to be recovered from the final
joint state. For two inputs \(z,z'\) in the compatible domain, equality of the
registers is equivalent to
\(\Omega_M(Uz)=\Omega_M(Uz')\). Hence, \(Uz\) and \(Uz'\) belong to the
same fiber of the composite map \(q_M\circ\operatorname{pr}_M\). This
composition includes both the loss due to projection onto the apparatus and
that due to evaluation of its state. Any subsequent erasure introduces an
additional map and must be analyzed separately.
\end{proof}

\begin{figure}[htbp]
\centering
\resizebox{\textwidth}{!}{%
\begin{tikzpicture}[node distance=13mm and 15mm,>=Latex,
box/.style={draw=hmtblue,rounded corners,fill=hmtlight,align=center,minimum height=9mm,minimum width=29mm}]
\node[box] (s) {complete state\\[-1mm]$x$};
\node[box,right=of s] (u) {acoplamiento\\[-1mm]$U(x,m_0)$};
\node[box,right=of u] (c) {restricción\\[-1mm]$\mathcal C$};
\node[box,right=of c] (q) {evaluación\\[-1mm]$q_M$};
\node[box,right=of q] (y) {observable\\[-1mm]$y$};
\draw[->,thick] (s)--(u); \draw[->,thick] (u)--(c);
\draw[->,thick] (c)--(q); \draw[->,thick] (q)--(y);
\end{tikzpicture}
}
\caption{Composición de una medición. La aplicación de evaluación actúa
después del acoplamiento y de la restricción de compatibilidad; el estado
conjunto permanece disponible para la transformación inversa.}
\label{p12-83-c16-s1:fig:medicion}
\end{figure}

\subsection{Principio de Ambivalencia: involución bidireccional de los lectores π/φ² y φ/π²}

Sean
\[
 q_A:\mathcal X_{\rm TPK}\to\mathcal X_A,
 \qquad
 q_V:\mathcal X_{\rm TPK}\to\mathcal X_V
\]
two quotient maps respectively associated with the observables of
aperture and retention, and let
\(T:\mathcal X_{\rm TPK}\to\mathcal M\) be the transported memory
variable.

\begin{principle}[Ambivalence]
The complete state precedes its observable projections. The equivalence
relations induced by \(q_A\) and \(q_V\) are not assumed to be identical, and
compatibility between them is recorded by \(T\). The joint map
\[
 (q_A,q_V,T):\mathcal X_{\rm TPK}
 \longrightarrow\mathcal X_A\times\mathcal X_V\times\mathcal M
\]
is adopted as the separation condition for physically distinguishable
states, modulo the declared gauge group.
\end{principle}

\begin{proposition}[Reconstruction under the separation condition]
If the preceding joint map is injective modulo gauge, each complete-state
class is uniquely determined by its two observables and its memory
variable.
\end{proposition}
\begin{proof}
Two states with the same triple of images belong to the same gauge class
by the injectivity hypothesis. Each fiber of the joint map
therefore contains at most one class.
\end{proof}

The observable \(q_A\) records the accessible prolongations, and \(q_V\) the
retained component. A material interpretation of retention requires a
mass observable that factors through \(q_V\); it does not follow solely from the
separation principle.

\subsection{Relational determination and temporal orientation}

Let \([\gamma]_{\mathcal G}\) be the class of a trajectory in the state groupoid, and
let \(\beta_\gamma\) be its terminal boundary datum.
\begin{principle}[Relational determination]
An admissible local observable has the form
\[
 \mathcal O_{\rm loc}(x;\gamma)
 =F\bigl(x,[\gamma]_{\mathcal G},\beta_\gamma\bigr)
\]
and it is invariant under trajectory equivalences that preserve the
ordered record. Consequently, the visible point coordinate does not by itself
determine the observable when the evaluation has identified different
transport classes.
\end{principle}

The formulation is relational in a precise sense: the argument of
\(\mathcal O_{\rm loc}\) contains the state, the global transport class, and
the boundary. This is the domain on which the theory may subsequently be compared
with the tradition associated with Mach's principle. It is not presupposed that
every inertial property has been derived; the type of the map that
must realize it is specified.

If the update on \(\mathcal X_{\rm TPK}\) is bijective and the register
preserves the coordinates required by its inverse, the complete dynamics is
reversible. An observable temporal orientation appears upon composing it with a
noninjective reader, reinitializing a boundary datum, or selecting a single
prolongation among several. Irreversibility is therefore localized in a
defined morphism of information loss or selection, not in the mere
ordering of labels.

\subsection{Relation with classical variational principles}

Hamilton's principle requires the first variation of an action to vanish
along the physical trajectory; it does not assert, in general, that this trajectory is a
global minimum. The discrete interaction functional constructed in the
\cref{chap:pantalla}, by contrast, orders a finite family of routes by the
number of changes in orientation and arithmetic law. The two principles
coincide only after constructing a limit in which the discrete cost
induces an action and its first variation reproduces the equations of
motion. This distinction precludes replacing a future correspondence with
an identity of names.

\subsection{Conditions for a prospective evaluation}

A prospective evaluation is mathematically determined when, before
its output is consulted, the state space, coupling,
quotient maps, memory variable, compatibility rule, and
comparison criterion have been fixed. If several sections survive, the output is the complete
family. The assignment of frequencies then requires a probability
measure or an explicit additional dynamics.


\subsection{Quotient of modes and areal-volumetric correspondence}
\label{p12-83-c16-s2:sec:cociente-modos-fav}

The calendar sign, mode label, and geometric sheet are three typed
objects. In particular, the phase sign is not identified with the
angular observable on digits. For the geometric correspondence, we fix the
typed HMT reader
\[
 \mathfrak g_{\rm av}(\Sigma)=\mathscr H_{\rm ar},
 \qquad
 \mathfrak g_{\rm av}(\Pi)=\mathscr H_{\rm vol}.
\]
The first assignment expresses the additive or boundary reading; the second, the
multiplicative or interior reading. The map \(\mathfrak g_{\rm av}\) is
part of the sheet typing. The following result proves that, once
that typing is fixed, TPK dynamics forces the incidence between them.

\begin{theorem}[Necessary descent of leaf incidence]
\label{p12-83-c16-s2:thm:descenso-necesario-fav}
Let \(m_r\) be the mode obtained after applying phase \(r\) of the
primitive block \((+1,-1,0)\), with the rule
\[
 +1:\ m\mapsto\Sigma,\qquad
 -1:\ m\mapsto\Pi,\qquad
 0:\ m\mapsto m.
\]
The cyclic orbit of modes is \((\Sigma,\Pi,\Pi)\). If
\[
 N_3(i,j)
 :=
 \#\{r\in\mathbb Z/3\mathbb Z:(m_r,m_{r+1})=(i,j)\},
\]
then, in the order \((\Sigma,\Pi)\),
\[
 \boxed{
 N_3=
 \begin{pmatrix}0&1\\1&1\end{pmatrix}.}
\]
By repetition of the calendar,
\[
 N_9=3N_3,\qquad N_{27}=9N_3,\qquad N_{108}=36N_3.
\]
After applying \(\mathfrak g_{\rm av}\), the primitive incidence
matrix is, in the order
\((\mathscr H_{\rm ar},\mathscr H_{\rm vol})\),
\[
 \boxed{
 F_{\rm av}=
 \begin{pmatrix}0&1\\1&1\end{pmatrix}.}
\]
Thus, the two cross transitions, the unique volumetric self-retention, the
absence of areal self-retention and the unit primitive multiplicities
are consequences of the TPK mode quotient; they are not four added
postulates.
\end{theorem}

\begin{proof}
The three consecutive pairs, with cyclic closure, are
\[
 \Sigma\longrightarrow\Pi,\qquad
 \Pi\longrightarrow\Pi,\qquad
 \Pi\longrightarrow\Sigma.
\]
Each appears once, and the pair \(\Sigma\to\Sigma\) does not appear. This determines the four entries of \(N_3\). The calendars of lengths \(9,27,108\) respectively contain \(3,9,36\) copies of the primitive block, so their incidence matrices are the indicated multiples. The greatest common divisor of the nonzero entries of each long matrix removes only that global repetition and returns \(N_3\). Finally, \(\mathfrak g_{\rm av}\) transports \(\Sigma,\Pi\) to the areal and volumetric sheets without altering the incidence.
\end{proof}

The real hyperbolic realization that uses this matrix and proves its
Lorentz invariance is formulated in
\cref{sec:tres-generadores-concretos-trit}; there \(F_{\rm av}\) is used
as an antecedent, without repeating its derivation.

This theorem must be distinguished from a stronger Markov assertion.
Forgetting the phase is not a strong Markov quotient of
\(\mathcal K_{\rm ph}\): the visible mode does not determine which of the three
operators \(I,E_+,E_\times\) acts at the next elementary iteration. What descends
without an additional hypothesis is the edge multigraph of a TPK\@ block.

There is, however, a memory-one completion uniquely determined by that descent. Let \(X_{\rm av}\) be the smallest subshift of finite type over \(\{\mathscr H_{\rm ar},\mathscr H_{\rm vol}\}\) that contains all the words obtained by forgetting the phase and retains only consecutive pairs. Every one-step subshift with that property must admit exactly the three observed pairs; the smallest such subshift excludes the sole unobserved pair. Its adjacency matrix is therefore \(F_{\rm av}\).

\begin{corollary}[Parry Measure of the One-Memory Wrapper]
\label{p12-83-c16-s2:cor:parry-envolvente-fav}
The subshift \(X_{\rm av}\) is primitive, has topological entropy
\(\log\varphi\), and possesses a unique measure of maximal entropy. Its stationary
vertex weights satisfy
\[
 \boxed{
 \frac{\mu_{\rm ar}}{\mu_{\rm vol}}=\varphi^{-2}.}
\]
\end{corollary}

\begin{proof}
The Perron root of \(F_{\rm av}\) is \(\varphi\), and a positive Perron
vector is \((1,\varphi)\). Since the matrix is symmetric, the Parry vertex
weights are proportional to the product of the left and right vectors,
that is, to \((1,\varphi^2)\).
\end{proof}

The Parry measure belongs to the memory-one path envelope. It is not
the image of \(\tau_{468}\), nor the temporal frequency of the periodic orbit
of modes. The latter is
\[
 \nu_{\rm cyc}(\mathscr H_{\rm ar})=\frac13,\qquad
 \nu_{\rm cyc}(\mathscr H_{\rm vol})=\frac23,
\]
whereas \(\tau_{468}\) is uniform over emissions, and the Parry measure
weights renewal histories of arbitrary length. The three measures
answer distinct questions. The irrational ratio
\(\varphi^{-2}\) is thus derived canonically from the path envelope,
without attributing it to a quotient of finite cardinalities.

\subsubsection{From chronological counting to the return marked by \texorpdfstring{\(\pi\)}{π}}
\label{p12-83-c16-s2:sec:retorno-marcado-pi-phi2}

The two formulas
\[
 \left(\frac13,\frac23\right)
 \qquad\text{and}\qquad
 \frac{\mu_{\rm ar}}{\mu_{\rm vol}}=\varphi^{-2}
\]
are not approximations of one another. The first counts the three instants of the
primitive calendar \((\Sigma,\Pi,\Pi)\); the second weights all the
admissible histories of the memory-one envelope. Hence the chronological
frequency is rational, and the Parry ratio is irrational. In what follows,
we write
\[
 \frac{\pi_{\rm HMT}}{\varphi^2}
 =\pi_{\rm HMT}\varphi^{-2};
\]
this is not division by \(\varphi^{-2}\).

The appearance of \(\varphi^{-2}\) can be seen directly in the linear
dynamics. If \(F_n\) denotes the \(n\)-th Fibonacci number,
\[
 F_{\rm av}^{\,n}
 =
 \begin{pmatrix}
  F_{n-1}&F_n\\
  F_n&F_{n+1}
 \end{pmatrix}
 \qquad(n\geq1).
\]
The eigenvalues of \(F_{\rm av}\) are \(\varphi\) and \(-\varphi^{-1}\). The contracted branch reverses orientation in one iteration. Upon passing to the quadratic memory, namely \(\operatorname{Sym}^2(\mathbb R^2)\), that reversal is recorded before the sign disappears. In the basis \((x^2,2xy,y^2)\), with the images of the basis vectors represented by rows, the induced operator is
\[
 \operatorname{Sym}^2(F_{\rm av})=
 \begin{pmatrix}
  0&0&1\\
  0&1&2\\
  1&1&1
 \end{pmatrix},
 \qquad
 \chi(t)=(t+1)(t^2-3t+1).
\]
Its three eigenvalues are
\[
 \varphi^2,\qquad -1,\qquad \varphi^{-2}.
\]
Thus, \(\varphi^{-2}\) is the positive stable mode of second-order
memory. This step explains why the golden ratio appears here as a return
law and not as an added statistical mean.

The closure \(\pi\) belongs to another level of the same HMT state: it is the
coinductive output of the regional reader already constructed, not a new eigenvalue of
\(F_{\rm av}\). Let
\[
 P_{\rm ar}=
 \begin{pmatrix}1&0\\0&0\end{pmatrix},
 \qquad
 P_{\rm vol}=
 \begin{pmatrix}0&0\\0&1\end{pmatrix}.
\]
Among the positive operators that preserve both sheets, there is exactly
one whose volumetric normalization is one and whose areal mark is the generated
closure \(\pi_{\rm HMT}\):
\[
 D_\pi=\pi_{\rm HMT}P_{\rm ar}+P_{\rm vol}.
\]
Uniqueness is relative to these three explicit clauses: positivity,
diagonality with respect to the leaves, and the two normalizations. With
\(e_{\rm ar}=(1,0)^T\) and \(e_{\rm vol}=(0,1)^T\), the marked return of
length \(n\) is
\[
 \Theta_n
 :=
 \frac{\langle e_{\rm ar},
 D_\pi F_{\rm av}^{\,n}e_{\rm ar}\rangle}
 {\langle e_{\rm vol},
 F_{\rm av}^{\,n}e_{\rm vol}\rangle}
 =
 \pi_{\rm HMT}\frac{F_{n-1}}{F_{n+1}}.
\]
Therefore,
\[
 \boxed{\displaystyle
 \lim_{n\to\infty}\Theta_n
 =\frac{\pi_{\rm HMT}}{\varphi^2}.}
\]
The mark \(\pi_{\rm HMT}\) is applied only once, at the boundary return;
it is not reinjected at every transition. Finite combinatorics produces the
stable mode \(\varphi^{-2}\), the coinductive reader produces the closure
\(\pi_{\rm HMT}\), and \(D_\pi\) composes both data without reversing their
causal order.

The native length \(108\) provides an exact witness to that convergence:
\[
 (F_{\rm av}^{108})_{\rm ar,ar}
 =F_{107}=10284720757613717413913,
\]
\[
 (F_{\rm av}^{108})_{\rm vol,vol}
 =F_{109}=26925748508234281076009.
\]
Consequently,
\[
 \left|
 \frac{F_{107}}{F_{109}}-\varphi^{-2}
 \right|
 =
 6.1684979916614407936\ldots\times10^{-46},
\]
and, after the only closing marker,
\[
 \left|
 \pi_{\rm HMT}\frac{F_{107}}{F_{109}}
 -\frac{\pi_{\rm HMT}}{\varphi^2}
 \right|
 =
 1.9378907974286976068\ldots\times10^{-45}.
\]
These errors measure the finite approximation of the return; they are not adjustments of
\(\pi\) or of \(\varphi\).

\paragraph{Cylindrical reader of the ratio.}
Let \(I_n^\pi=[\underline\pi_n,\overline\pi_n]\) and
\(I_n^\varphi=[\underline\varphi_n,\overline\varphi_n]\) be the positive,
nested cylinders produced by the two HMT readers\@. The interval
operation
\[
 \mathcal Q(I,J)
 :=
 \left[
 \frac{\inf I}{(\sup J)^2},
 \frac{\sup I}{(\inf J)^2}
 \right]
\]
is the smallest interval containing \(x/y^2\) for every
\((x,y)\in I\times J\). It preserves nesting: if the input cylinders
are refined, then \(\mathcal Q(I_n^\pi,I_n^\varphi)\) is also refined. Since their
diameters tend to zero and the limit of \(I_n^\varphi\) is positive,
\[
 \bigcap_n\mathcal Q(I_n^\pi,I_n^\varphi)
 =
 \left\{\frac{\pi_{\rm HMT}}{\varphi_{\rm HMT}^2}\right\}.
\]
With the twelve already published base-thousand blocks, the quotient interval forces
thirty-five decimal digits:
\[
 \frac{\pi_{\rm HMT}}{\varphi_{\rm HMT}^2}
 =
 1.19998161486432666111577775331680692\ldots.
\]
This construction preserves the provenance of each digit: the numerator
comes from the closure cylinder and the denominator from the autoscale.

\paragraph{Algebraic separation between areal reading, volumetric reading and their respective codomains}
The complete ratio cannot arise from the rational algebra of
\(F_{\rm av}\) alone. Indeed, its rational spectral functions belong to
\(\mathbb Q(\sqrt5)\); \(\varphi^{-2}\) is found there, but not the
transcendental number \(\pi\). Consequently, any derivation purporting to
obtain \(\pi/\varphi^2\) solely from the matrix \(2\times2\) would be
ill typed. The transcendental datum must derive from the coinductive closure
reader, or from an equivalent boundary cocycle proved
independently. This negative control fixes the division of labor between
the finite correspondence and the number reader.

\subsubsection{Exchange involution of the quotient}
\label{p12-83-c16-s2:sec:espejo-pi-phi-electron-hmt}

The correspondence admits two exact orientations of the same law. If
\[
 \mathscr R(x,y)=\frac{x}{y^2},
 \qquad
 \mathsf J(x,y)=(y,x),
\]
then \(\mathsf J^2=I\) and
\[
 \mathscr R(\pi,\varphi)=\frac{\pi}{\varphi^2},
 \qquad
 (\mathscr R\circ\mathsf J)(\pi,\varphi)
 =\frac{\varphi}{\pi^2}.
\]
The identities are also verified.
\[
 \mathscr R(\pi,\varphi)\mathscr R(\varphi,\pi)
 =\frac1{\pi\varphi},
 \qquad
 \frac{\mathscr R(\pi,\varphi)}
 {\mathscr R(\varphi,\pi)}
 =\left(\frac{\pi}{\varphi}\right)^3.
\]
Writing
\(\mathcal B_{\rm av}^{(\pi)}=\mathscr R(\pi,\varphi)\) also yields
the exact reparametrization
\[
 \frac{\varphi}{\pi^2}
 =
 \frac1{\varphi^3
 \bigl(\mathcal B_{\rm av}^{(\pi)}\bigr)^2},
 \qquad
 e^{\varphi/\pi^2}
 =
 \exp\!\left[
 \frac1{\varphi^3
 \bigl(\mathcal B_{\rm av}^{(\pi)}\bigr)^2}
 \right].
\]
Thus, the change of orientation between closure and growth transports the
boundary ratio to its exchanged reading:
\[
 \frac{\pi}{\varphi^2}
 \xleftrightarrow{\ \mathsf J\ }
 \frac{\varphi}{\pi^2}.
\]
The first orientation reads the marked areal closure over quadratic
volumetric memory; the second reads interior growth over
quadratic closure. This correspondence is a theorem of the HMT reader. Dimensional Mechanics subsequently applies the second orientation within the
electronic construction, without converting the electron into an antecedent of
\(\pi\), \(\varphi\), or of the correspondence.

\subsubsection{Ambivalence Principle of the areal and volumetric readers}
\label{p12-83-c16-s2:sec:ambivalencia-grav-inercial-pi-phi2}

The Ambivalence Principle asserts that the areal reading and the
volumetric reading are two coordinated, but not interchangeable, functionals of one
and the same state with memory. Their coincidence on the diagonal is
\[
 \mathcal L_{\rm ar}^{\rm diag}(x)
 =\mathcal L_{\rm vol}^{\rm diag}(x).
\]
The diagonal equality does not erase the ambivalence outside of it.

On the areal--volumetric boundary, both readers factor through the same
positive memory amplitude \(\mathcal M(x)\):
\[
 \mathcal L_{\rm vol}(x)=\mu_{\rm vol}\mathcal M(x),
 \qquad
 \mathcal L_{\rm ar}(x)=
 \pi_{\rm HMT}\mu_{\rm ar}\mathcal M(x).
\]
Then the Parry theorem and the return mark give
\[
 \boxed{\displaystyle
 \frac{\mathcal L_{\rm ar}(x)}
 {\mathcal L_{\rm vol}(x)}
 =\frac{\pi_{\rm HMT}}{\varphi^2}.}
\]
With volumetric normalization, the relative difference of readers is
\[
 \frac{\mathcal L_{\rm ar}-\mathcal L_{\rm vol}}
 {\mathcal L_{\rm vol}}
 =
 \frac{\pi_{\rm HMT}}{\varphi^2}-1
 =
 0.1999816148643266611\ldots.
\]
The physical reading of this correspondence is undertaken later, in the
Cartan--Holst chapter, after constructing holonomy, co-tetrad, connection, torsion,
and curvature. Here the result is entirely HMT: an exact ratio between
two readers of the same sheet memory.

They should be distinguished from the enriched extractor subsequently used in
Dimensional Mechanics. There a complete history preserves the action count
\(n_A\), the twelve oriented events \(e_0,\ldots,e_{11}\) and the torsional
cochain at step level; these yield
\[
 s_C=\sum_{i=0}^{5}(e_i+e_{i+6}),
 \qquad
 W_\pm=6n_A\pm s_C.
\]
The channels already generated by HMT then satisfy the exact identity
\[
 \exp\!\left(n_AA_{\rm rad}+\frac{s_C}{6}C^\ast_{\rm rad}\right)
 =\left(q_+^{-W_+}q_-^{-W_-}\right)^{1/12}.
\]
Therefore, the route--angular-character passage is not absent: it is realized on
the enriched dodecaphase record. What this phase calculation proves by itself
are \(\Omega_{\rm sw},S,G\); it neither replaces them nor permits reconstruction of the
complete record after it has been projected.

The uniform measure on the \(104\,976\) presentations descends to atoms
of masses \(1/729\), \(1/243\) and \(1/54\), but not to \(1/468\). In the
fivefold microfiber of \(\pi\), the two probabilities are
\[
 \tau_{468}(\mathscr R_\pi)=\frac5{468},
 \qquad
 F_\#\mu_{\rm seed}(\mathscr R_\pi)=\frac7{729}.
\]
Both the averaged lift and the phase lift realize the first: they do not
push forward the uniform measure on seeds, but assign the same total
mass to each observable class. The phase lift also preserves
microscopic stillness at neutral steps.

To relate phase elevation to the enriched histories of
depth six, a kernel \(\mathcal K_6\) on
\(\mathcal H_6^{\rm cat}\times C_9\) is required. It descends to
\(\widehat{\mathcal K}_{\rm ph}\) through
\(\rho_6\times\mathrm{id}_{C_9}\) if and only if it satisfies, at each phase, the
strong lumpability condition
\[
 \sum_{\rho_6(g)=y}\mathcal K_6((h,r),(g,r+1))
 =L_{\tau(r)}(\rho_6(h),y),
\]
independently of \(h\) within each fiber of \(\rho_6\). In that case,
composition with \(F\) projects to \(\mathcal K_{\rm ph}\). Constructing
a family \((\mathcal K_m)_{m\geq6}\) that satisfies this identity and
commutes with all truncations is the remaining step toward a
transfer of the inverse limit.


\section{Areal--Volumetric Correspondence and Ratio $\pi /\varphi ^2$}
\label{p12-83-c16-s3:sec:bisagra-hojas-pi-phi2}

The preceding section constructed the dodecaphase involution, its lift to persistent sections, the Möbius band relative to a negative holonomy, and the biangular chart \((A,C^*)\leftrightarrow(q_+,q_-)\). It remains to state precisely how the returns of the areal and volumetric sheets are counted and why the golden ratio appears squared.

\subsection{Classes, supports, and leaves are distinct objects}

The nonadic classes of motion are
\[
 C_1=\{1,4,7\},\qquad C_2=\{2,5,8\},\qquad C_0=\{3,6,9\}.
\]
The supports employed by the angular observable are
\[
 \mathcal D_{\rm act}=\{1,3,5,7\},\qquad
 \mathcal D_{\rm evt}=\{2,4,6,8\},\qquad
 \mathcal D_9=\{9\}.
\]
Its exact intersection is
\[
 B_{\rm na}=\bigl(|C_i\cap D_j|\bigr)_{i,j}
 =\begin{pmatrix}2&1&0\\1&2&0\\1&1&1\end{pmatrix}.
\]

\begin{theorem}[Triadic incidence]
\label{p12-83-c16-s3:thm:incidencia-triadica-corta}
The following identities hold
\[
 \det B_{\rm na}=3,
 \qquad
 \operatorname{SNF}(B_{\rm na})=\operatorname{diag}(1,1,3),
\]
\[
 \operatorname{im}B_{\rm na}
 =\{(u,v,w)\in\mathbb Z^3:u+v\equiv0\pmod3\}.
\]
The simultaneous interchange
\(C_1\leftrightarrow C_2\),
\(\mathcal D_{\rm act}\leftrightarrow\mathcal D_{\rm evt}\) leaves the
incidence fixed.
\end{theorem}

\Needspace{5\baselineskip}
The observable of the record
\[
 j_2(d)=
 \begin{cases}
 -1,&d\in\mathcal D_{\rm act},\\
 +1,&d\in\mathcal D_{\rm evt},\\
 0,&d\in\mathcal D_9
 \end{cases}
\]
is composed with the active geometric reader:
\[
 -1\mapsto\mathscr H_{\rm ar},\qquad
 +1\mapsto\mathscr H_{\rm vol},\qquad
 0\mapsto\mathscr H_{\rm tor}.
\]
Therefore,
\[
 \mathfrak g_\triangle(\mathcal D_{\rm act})=\mathscr H_{\rm ar},\qquad
 \mathfrak g_\triangle(\mathcal D_{\rm evt})=\mathscr H_{\rm vol},\qquad
 \mathfrak g_\triangle(\mathcal D_9)=\mathscr H_{\rm tor}.
\]
The matrix proves the crossing between classes and supports. The reader
\(\mathfrak g_\triangle\) declares the passage from supports to sheets. Omitting this
second map would identify the cardinalities of distinct objects.

\subsection{Spectral opening of the band through the conjugated angular phase}

In the active chart,
\[
 q_+=e^{-A_{\rm rad}-C^*_{\rm rad}},\qquad
 q_-=e^{-A_{\rm rad}+C^*_{\rm rad}},
\]
and the exact inverse is
\[
 A_{\rm rad}=-\frac12\log(q_+q_-),
 \qquad
 C^*_{\rm rad}=\frac12\log\frac{q_-}{q_+}.
\]
Thus, \(A\) is the common mode and \(C^*\) the differential mode. At each
tower height,
\[
 q_-^n-q_+^n
 =2e^{-nA_{\rm rad}}\sinh(nC^*_{\rm rad}).
\]
The negative holonomy constructed in
Section~\ref{sec:persistencia-ruta-mobius-familias} fixes the Möbius
gluing; the last identity proves that \(C^*\) measures the oriented opening
of the two conjugate sheets. In the physical representation, those two
orientations are read as particle and antiparticle.

In the boundary reader, one turn of the base is parameterized by \(2\pi\) and
ends in the opposite orientation. The squared loop, parameterized by
\(4\pi\), restores the orientation. Therefore, relative to the declared lift and
geometric reader, \(2\pi\) is projected areal closure and
\(4\pi\) is volumetric closure of the oriented cover. This construction does not
conflate that double return with a spinorial representation; the
representation, when introduced, is a subsequent functor.

\subsection{TPK descent, path envelope, and loop counting}

The TPK calendar and the angular observable \(j_2\) are distinct maps.
The former updates the arithmetic mode \(m\in\{\Sigma,\Pi\}\); the latter classifies
digit supports. The hinge uses the typed HMT reader
\[
 \eta_{\rm av}(\Sigma)=\mathscr H_{\rm ar},
 \qquad
 \eta_{\rm av}(\Pi)=\mathscr H_{\rm vol}.
\]
Thus, \(\Sigma\) is the additive or boundary mode and \(\Pi\) the
multiplicative or interior mode. This statement does not identify
\(\tau\) with \(j_2\).

In the primitive phase block
\[
 \tau_3=(+1,-1,0),
\]
TPK updates the mode via
\[
 +1:\ m\mapsto\Sigma,\qquad
 -1:\ m\mapsto\Pi,\qquad
 0:\ m\mapsto m.
\]
The cyclic orbit is \((\Sigma,\Pi,\Pi)\). Define its edge-incidence
matrix by
\[
 N_3(i,j)=
 \#\{r\in\mathbb Z/3\mathbb Z:(m_r,m_{r+1})=(i,j)\}.
\]

\begin{theorem}[Necessary Descent of the Four Rules]
\label{p12-83-c16-s3:thm:descenso-necesario-fav-corta}
In the order \((\Sigma,\Pi)\),
\[
 N_3=
 \begin{pmatrix}0&1\\1&1\end{pmatrix}.
\]
Moreover,
\[
 N_9=3N_3,\qquad
 N_{27}=9N_3,\qquad
 N_{108}=36N_3.
\]
Therefore, after removing the common repetition factor and applying
\(\eta_{\rm av}\), the TPK mode quotient necessarily induces
\[
 \boxed{
 F_{\rm av}=S+P_{\rm vol}
 =
 \begin{pmatrix}0&1\\1&1\end{pmatrix}
 }
\]
where
\[
 S=\begin{pmatrix}0&1\\1&0\end{pmatrix},
 \qquad
 P_{\rm vol}=\begin{pmatrix}0&0\\0&1\end{pmatrix}.
\]
The two crossed transitions, the unique volumetric self-retention, the
absence of areal self-retention, and the unit primitive multiplicities
are consequences of TPK, not added premises.
\end{theorem}

\begin{proof}
The three consecutive pairs of the cycle are
\[
 \Sigma\longrightarrow\Pi,\qquad
 \Pi\longrightarrow\Pi,\qquad
 \Pi\longrightarrow\Sigma.
\]
Each one appears once and \(\Sigma\to\Sigma\) does not appear. This determines
the four entries of \(N_3\). The length calendars
\(9,27,108\) contain \(3,9,36\) primitive blocks, respectively;
therefrom come the three multiples. Reduction by the greatest common divisor of the
nonzero entries eliminates only the repetition of the block. The reader
\(\eta_{\rm av}\) changes the type of the vertices, not their incidence.
\end{proof}

The preceding theorem is an exact descent of edges. It does not turn phase
forgetting into a strong Markov quotient: the same visible emission may
be updated by \(I,E_+\) or \(E_\times\), depending on the hidden phase. From
this fact follows a necessary distinction among three measures:
\[
 \tau_{468},\qquad
 \nu_{\rm cyc},\qquad
 \mu_{\rm Parry}.
\]
The former is uniform over the \(468\) emissions and stationary under
phase transfer. The latter is the frequency of the periodic orbit,
\[
 \nu_{\rm cyc}(\mathscr H_{\rm ar})=\frac13,\qquad
 \nu_{\rm cyc}(\mathscr H_{\rm vol})=\frac23.
\]
Neither of them can directly produce the irrational quotient
\(\varphi^{-2}\).

The object that does produce it is determined by a universal property.
Let \(X_{\rm av}\) be the smallest memory-one subshift of finite type that
contains the words obtained by forgetting the phase and preserves all their
consecutive pairs. Every one-step system containing those words must
admit
\[
 {\rm ar}\to{\rm vol},\qquad
 {\rm vol}\to{\rm vol},\qquad
 {\rm vol}\to{\rm ar}.
\]
The smaller one excludes the only absent pair, \({\rm ar}\to{\rm ar}\). Consequently,
\(X_{\rm av}\) has adjacency matrix \(F_{\rm av}\).
This is the minimal one-memory symbolic completion of the TPK quotient; it is
not a new microscopic dynamics.

\begin{theorem}[Fibonacci Law and Parry Measure]
\label{p12-83-c16-s3:thm:ley-fibonacci-retornos-corta}
The operator \(F_{\rm av}\) is primitive and
\[
 F_{\rm av}^{\,n}
 =\begin{pmatrix}F_{n-1}&F_n\\F_n&F_{n+1}\end{pmatrix}
 \qquad(n\geq1).
\]
The subshift \(X_{\rm av}\) has topological entropy
\(\log\varphi\) and a unique invariant measure of maximal entropy. For
this measure,
\[
 \frac{(F_{\rm av}^{\,n})_{\rm ar,ar}}
 {(F_{\rm av}^{\,n})_{\rm vol,vol}}
 =\frac{F_{n-1}}{F_{n+1}}
 \longrightarrow\varphi^{-2},
 \qquad
 \boxed{\frac{\mu_{\rm ar}}{\mu_{\rm vol}}=\varphi^{-2}.}
\]
\end{theorem}

\begin{proof}
The square of \(F_{\rm av}\) is strictly positive. The formula for
powers is obtained by induction through
\(F_{n+1}=F_n+F_{n-1}\). The diagonal entries count closed loops.
The Perron root is \(\varphi\) and a positive Perron vector is
\((1,\varphi)\). Since the matrix is symmetric, the vertex weights of
Parry are proportional to the product of the left and right vectors,
that is, to \((1,\varphi^2)\).
\end{proof}

The factor \(\varphi^{-1}\) is the Perron amplitude ratio of a
traversal; \(\varphi^{-2}\) is the weight ratio of a complete return.
The two exponents are thereby ordered without introducing two independent
constants.

\subsection{Spectral Origin of the Quadratic Exponent of Return}

The appearance of the square admits a second, purely spectral proof that
does not depend on interpreting a weighted average. The characteristic
polynomial of \(F_{\rm av}\) is
\[
 \chi_{F_{\rm av}}(t)=t^2-t-1,
\]
so that its two eigenvalues are
\[
 \lambda_+=\varphi,\qquad \lambda_-=-\varphi^{-1}.
\]
The sign of \(\lambda_-\) records reversal of orientation in a traversal of
the stable sheet. Upon lifting the action to quadratic memory
\(\operatorname{Sym}^2(\mathbb R^2)\), the eigenvalues are the pairwise products
\[
 \lambda_+^2=\varphi^2,\qquad
 \lambda_+\lambda_-=-1,\qquad
 \lambda_-^2=\varphi^{-2}.
\]
In the basis
\((e_{\rm ar}^2,e_{\rm ar}\odot e_{\rm vol},e_{\rm vol}^2)\), writing the coordinates of the images of the three basis vectors
by columns, the matrix is
\[
 \operatorname{Sym}^2(F_{\rm av})=
 \begin{pmatrix}
 0&0&1\\
 0&1&2\\
 1&1&1
 \end{pmatrix},
\]
and satisfies
\[
 \chi_{\operatorname{Sym}^2(F_{\rm av})}(t)
 =t^3-2t^2-2t+1
 =(t+1)(t^2-3t+1).
\]
Therefore, the stable mode changes orientation with
\(-\varphi^{-1}\) in one traversal and becomes the positive weight
\(\varphi^{-2}\) when read as area memory or as a
quadratic return. This is the structural reason for the exponent two.

\subsection{Boundary closure and quotient \texorpdfstring{$\pi/\varphi^2$}{pi/phi2}}

The coinductive reader of \(\pi\) precedes this construction. We define the
closure band by composition, without using \(\pi\) to select the
transporter:
\[
 \mathcal B_{\rm av}^{(\pi)}(X_\infty)
 :=\Pi_{\rm HMT}(X_\infty)
 \frac{\mu_{\rm ar}}{\mu_{\rm vol}}.
\]

\paragraph{Marked return \(\Theta_n\) and convergence of the areal–volumetric ratio}
Let the sheet projectors be
\[
 P_{\rm ar}=\begin{pmatrix}1&0\\0&0\end{pmatrix},
 \qquad
 P_{\rm vol}=\begin{pmatrix}0&0\\0&1\end{pmatrix}.
\]
The boundary marker is defined by
\[
 D_\pi:=\Pi_{\rm HMT}(X_\infty)P_{\rm ar}+P_{\rm vol}.
\]
It is the unique positive operator diagonal with respect to the sheet
decomposition that assigns unit weight to the volumetric sheet and the closure
\(\Pi_{\rm HMT}(X_\infty)\) to the areal sheet. Its canonicity is therefore
relative to three explicit clauses: sheet diagonality, volumetric normalization,
and areal marking by the already constructed closure reader.

If \(e_{\rm ar},e_{\rm vol}\) are the sheet vectors, the marked-return
functional of depth \(n\) is
\[
 \Theta_n(X_\infty):=
 \frac{\langle e_{\rm ar},
 D_\pi F_{\rm av}^{\,n}e_{\rm ar}\rangle}
 {\langle e_{\rm vol},
 F_{\rm av}^{\,n}e_{\rm vol}\rangle}.
\]
The previous power law gives, without approximation,
\[
 \boxed{
 \Theta_n(X_\infty)
 =\Pi_{\rm HMT}(X_\infty)\frac{F_{n-1}}{F_{n+1}}.}
\]
The closure \(\pi\) is applied only once, when marking the boundary return; it is not
reintroduced at each iteration of \(F_{\rm av}\).

\begin{corollary}[Areal--Volumetric Holographic Ratio]
\label{p12-83-c16-s3:cor:pi-phi2-corta}
Relative to the TPK mode quotient, its minimal memory-one envelope,
and the HMT closure of \(\pi\),
\[
 \boxed{\mathcal B_{\rm av}^{(\pi)}
 =\lim_{n\to\infty}\Theta_n
 =\frac{\pi_{\rm HMT}}{\varphi^2}.}
\]
\end{corollary}

\paragraph{Native Periodicity of \(108\) Phases and Closure of the Unitary Calendar}
At dodecaphasic depth \(108\), the two diagonal counts are
\[
 (F_{\rm av}^{108})_{\rm ar,ar}
 =F_{107}=10284720757613717413913,
\]
\[
 (F_{\rm av}^{108})_{\rm vol,vol}
 =F_{109}=26925748508234281076009.
\]
Consequently,
\[
 \frac{F_{107}}{F_{109}}
 =0.381966011250105151795413165634361882\ldots,
\]
y the error respecto of \(\varphi^{-2}\) is
\[
 \left|\frac{F_{107}}{F_{109}}-\varphi^{-2}\right|
 =6.1684979916614407936\ldots\times10^{-46}.
\]
After the boundary marking, the error relative to
\(\pi/\varphi^2\) is
\[
 \left|\pi\frac{F_{107}}{F_{109}}-\frac{\pi}{\varphi^2}\right|
 =1.9378907974286976068\ldots\times10^{-45}.
\]
Thus, \(108\) is not used to fit the ratio: it is a native depth of
the calendar that already approximates the limiting return with more than
forty-four stable fractional digits.

\paragraph{Coinductive intervalar reader: cylinder fitting and uniqueness of the limit}
For positive intervals
\[
 I=[a,b],\qquad J=[c,d],\qquad 0<c\leq d,
\]
we define
\[
 \mathcal Q(I,J):=
 \left[\frac{a}{d^2},\frac{b}{c^2}\right].
\]
The function \((x,y)\mapsto x/y^2\) is increasing in \(x\) and decreasing in
\(y>0\). Hence \(\mathcal Q(I,J)\) is exactly the smallest interval
containing all quotients \(x/y^2\) with
\((x,y)\in I\times J\); it is not a bound chosen by rounding.

Let \(I_n^\pi\) and \(I_n^\varphi\) be the nested positive cylinders
generated by HMT\@. Then
\[
 \mathcal Q(I_{n+1}^\pi,I_{n+1}^\varphi)
 \subseteq
 \mathcal Q(I_n^\pi,I_n^\varphi),
\]
and, as
\[
 \bigcap_n I_n^\pi=\{\pi_{\rm HMT}\},
 \qquad
 \bigcap_n I_n^\varphi=\{\varphi_{\rm HMT}\},
\]
we obtain
\[
 \bigcap_n\mathcal Q(I_n^\pi,I_n^\varphi)
 =\left\{\frac{\pi_{\rm HMT}}{\varphi_{\rm HMT}^2}\right\}.
\]
Each base-\(1000\) block is forced as soon as the quotient interval
lies entirely within a single triadic cylinder of the corresponding
depth.

In particular, the positive twelve-triad cylinders for \(\pi\) and
\(\varphi\) produce the quotient interval
\[
 \left[
 \frac{p^-}{(f^+)^2},
 \frac{p^+}{(f^-)^2}
 \right].
\]
The evaluation certifies thirty-five fractional digits:
\[
 \frac{\pi}{\varphi^2}
 =1.19998161486432666111577775331680692\ldots,
\]
y eleven complete blocks:
\[
 199|981|614|864|326|661|115|777|753|316|806.
\]

\paragraph{Algebraic Separation Theorem.}
The matrix \(F_{\rm av}\), its powers, and its rational spectral algebra
live in an algebraic extension of \(\mathbb Q\) contained in
\(\mathbb Q(\sqrt5)\). Consequently, no composition using only
rational algebraic operations on \(F_{\rm av}\) can produce the
transcendental factor \(\pi/\varphi^2\). The finite operator generates
\(\varphi^{-2}\); the coinductive boundary reader generates
\(\pi_{\rm HMT}\); the marked return composes both without conflating their
provenances. This negative control prevents \(F_{\rm av}\) from being presented as an
isolated generator of \(\pi\).

\subsection{Reader's Involution and Electronic Section}
\label{p12-83-c16-s3:sec:espejo-pi-phi-electron}

The relation just obtained appears alongside the electron in a manner more
precise than literal repetition. The active electron invariant of
Section~\ref{sec:persistencia-ruta-mobius-familias} contains
\[
 \mathfrak e_0=
 \frac{\sqrt3}{4}
 \left(e^{\varphi/\pi^2}-22\alpha_{\rm HMT}^{3}\right)
 (1+15\Delta_4).
\]
Therefore, the electron exponent is \(\varphi/\pi^2\), whereas the
sheet hinge produces \(\pi/\varphi^2\). The two scalars are neither
equal nor reciprocal. They form the two orientations of the same
bivariate law. Indeed, define
\[
 \mathscr R(x,y):=\frac{x}{y^2},
 \qquad
 \mathsf J(x,y):=(y,x).
\]
Then \(\mathsf J^2=I\) and
\[
 \boxed{
 \mathscr R(\pi,\varphi)=\frac{\pi}{\varphi^2},
 \qquad
 (\mathscr R\circ\mathsf J)(\pi,\varphi)
 =\frac{\varphi}{\pi^2}.}
\]
Exchange of the two readers carries the marked areal--volumetric return
exactly to the exponential seed of the electron. Moreover,
\[
 \mathscr R(\pi,\varphi)\mathscr R(\varphi,\pi)
 =\frac1{\pi\varphi},
 \qquad
 \frac{\mathscr R(\pi,\varphi)}
 {\mathscr R(\varphi,\pi)}
 =\left(\frac{\pi}{\varphi}\right)^3.
\]
If we denote by
\(\mathcal B_{\rm av}^{(\pi)}=\pi/\varphi^2\) the marked sheet
ratio, the same identity may be written without introducing an additional
constant:
\[
 \boxed{
 \frac{\varphi}{\pi^2}
 =
 \frac{1}
 {\varphi^3\bigl(\mathcal B_{\rm av}^{(\pi)}\bigr)^2},}
 \qquad
 e^{\varphi/\pi^2}
 =
 \exp\!\left[
 \frac{1}
 {\varphi^3\bigl(\mathcal B_{\rm av}^{(\pi)}\bigr)^2}
 \right].
\]
Numerically,
\[
 \begin{aligned}
 \frac{\varphi}{\pi^2}
 &=0.163941118913672383599752755745666982\ldots,\\
 e^{\varphi/\pi^2}
 &=1.17814494259630678081160095680888910\ldots.
 \end{aligned}
\]

This identity formalizes the bidirectional content without altering the force of the electron equation. The orientation \((\pi,\varphi)\) reads boundary closure through quadratic memory; the exchanged orientation \((\varphi,\pi)\) reads interior growth through quadratic closure and enters exponentiated in \(\mathfrak e_0\). The theorem proves that both expressions belong to the same exchange orbit. The coefficients \(22\) and \(15\) have an independent generator: the \cref{p12-83-c15-s2:thm:selector-persistencia-electron-rev7} obtains them, without consulting the electron mass, as
\[
 -22=-\bigl|\mathcal K_e\setminus\iota(\mathscr R_\pi)\bigr|,
 \qquad
 +15=\bigl|\mathscr R_\pi\times\mathbb F_3\bigr|.
\]
Its status is \textsc{demonstrated with explicit starting structure}. The
current involution transports that construction to the electronic orientation and
does not use the electron to obtain \(F_{\rm av}\), \(\pi\), or \(\varphi\). The
last form is an exact change of coordinates of the electronic seed, not
an alternative derivation of \(\mathfrak e_0\).

\subsection{Pentagonal Identity of Rogers--Ramanujan and Recognition of the Perron Invariant}

The Rogers--Ramanujan continued fraction provides an analytic realization
of the same invariant. With \(q_A=e^{-A_{\rm rad}}\), \(q_C=e^{-C^*_{\rm rad}}\), one obtains
\[
 R(q_A)R(q_C)-\varphi^{-2}\simeq-1.01864\times10^{-27},
\]
and with the channels of the split decomposition,
\[
 R(q_+)R(q_-)-\varphi^{-2}\simeq-1.19458\times10^{-21}.
\]
The differences are not zero. The exact statement is the limit.
\[
 q_A^{(m)},q_C^{(m)}\to1^-
 \quad\Longrightarrow\quad
 R(q_A^{(m)})R(q_C^{(m)})\to\varphi^{-2}.
\]
Rogers--Ramanujan thus recognizes the Perron invariant in the
pentadic reader; by itself, it does not select the angular coordinates \(A\), \(C^*\), or
\(F_{\rm av}\). The former finite kernel \(90/120\) preserves an approximate
residue \(7.81\times10^{-9}\) and is not used as an exact identity.

\subsection{Inversion of the inertial and gravitational readers and bidirectional reconstruction of the magnitude}

The Ambivalence Principle distinguishes the opening reader \(q_A\), the retention reader \(q_V\), and the transported memory. The preceding hinge makes it possible to formulate the physical intuition precisely without conflating it with the combinatorial theorem. Let \(\mathcal M(x)\) be a common memory character, and define, at the areal--volumetric interface,
\[
 \mathcal I_{\rm av}(x)
 :=\mu_{\rm vol}\mathcal M(x),
 \qquad
 \mathcal G_{\rm av}(x)
 :=\pi_{\rm HMT}\mu_{\rm ar}\mathcal M(x).
\]
Therefore, whenever \(\mathcal M(x)\ne0\),
\[
 \boxed{
 \frac{\mathcal G_{\rm av}(x)}{\mathcal I_{\rm av}(x)}
 =\frac{\pi_{\rm HMT}}{\varphi^2},}
 \qquad
 \frac{\mathcal G_{\rm av}-\mathcal I_{\rm av}}
 {\mathcal I_{\rm av}}
 =\frac{\pi_{\rm HMT}}{\varphi^2}-1.
\]
The first formula expresses the oriented ratio between the gravitational
and inertial readers; the second, their relative defect. The gnomonic L
carries the local Euclidean chart on which both readers descend to a
common normalization,
\[
 \frac{\mathcal G_{L}}{\mathcal I_{L}}=1.
\]

The physical reading is stratified. In the Euclidean chart of L, the
local torsional connection projects onto a unique normalization, and the
equivalence principle is realized as the unit quotient. In the enriched
state prior to that projection, areal opening and volumetric retention
remain separate: the Ambivalence Principle preserves its two
functionals and fixes the internal quotient \(\pi/\varphi^2\). Local equivalence
and boundary ambivalence are thus two successive readings of the
same antecedent, with declared domains and maps.

The physical realization adopts factorization of the inertial observable by
volumetric retention and of the gravitational observable by areal opening
with boundary closure \(\pi\); both share character \(\mathcal M\). The
internal mathematics determines the correspondence and its quotient, and the
Dimensional Mechanics functor publishes its observable charts. A rupture of either
factorization constitutes the explicit falsifier of this realization.

\paragraph{Cosmological boundary.}
The hypothesis that dark matter or dark energy are effective readings of nontorsional crossings is retained only as a hypothesis of physical prolongation. To turn it into a testable proposition, at least a map from sheet interfaces to a cosmological geometry, observables representing density and pressure, and a prospective comparison with data without using those data to fix the map must be constructed. The proved correspondence supplies the structural candidate; it does not yet constitute, by itself, an observational explanation of the dark sector.

\paragraph{Areal–volumetric hinge domain and necessary validity conditions}
The band, the areal--volumetric ratio, and the bidirectional reading belong to
the prior architecture. The incidence matrix, calendar descent,
minimal envelope, loop law, Parry measure, quadratic action,
and marked return provide its explicit formalization. The
interval reader and the depth controls \(108\) certify the
calculation. The incidences \(N_3,N_9,N_{27},N_{108}\), the matrix
identities, and the cylinders are exact internal identities; the Möbius theorem and the ratio
\(\pi/\varphi^2\) are
\textsc{proved from explicit starting structure}; the
particle--antiparticle and inertial--gravitational readings constitute subsequent physical
interfaces; the extension to the dark sector remains an
unproved physical hypothesis. The pair related by the involution
\(\mathscr R(\pi,\varphi)\leftrightarrow
\mathscr R(\varphi,\pi)\) exactly formalizes a prior architecture. The
coefficients \((-22,+15)\) are proved, by
\cref{p12-83-c15-s2:thm:selector-persistencia-electron-rev7}, from their explicit starting
structure.

A positive holonomy gives a cylinder. An areal--areal transition, the absence of a cross transition, or a TPK count different from \(N_9=3F_{\rm av}\) and \(N_{108}=36F_{\rm av}\) falsifies the descent. The identification of \(\mu_{\rm Parry}\) with \(\tau_{468}\) or with \(\nu_{\rm cyc}\) is refuted by their distinct values. A spectrum of \(\operatorname{Sym}^2(F_{\rm av})\) different from \(\{\varphi^2,-1,\varphi^{-2}\}\), a count \(108\) different from \((F_{107},F_{109})\), a nonnested quotient interval, or the iterative introduction of \(\pi\) within \(F_{\rm av}\) falsifies the new formalizations. A universal physical identification of the quotient without a map from sheets to observables lies outside the theorem. The delta of the four rules is closed; it does not reopen \(\pi\), \(\varphi\), \(C^*\), or the extension--route--ledger--character junction.
